\cotefile{29}
\scanpage{19}
de $Z_j$, on a $Z'_j\quad Z'_{i\deleted{j}}/N_{ij}$, où $Z'_i$ est l'image inverse de $Z'_{\unsure{\deleted{i}j|\deleted{j}i}}$ sur $Z_i$). Dans ce cas, le foncteur ``réalisation géométrique'' est bien défini, et induit surchaque $\mathrm{Rev}^{\underline{R}_i}(S)$ le foncteur réalisation géométrique déjà explicité dans l'exemple \unsure{\deleted{1}2|\deleted{2}1}. \deleted{Bien sûr, les conditions} Notons que les \unsure{cri\deleted{t}è\deleted{è}tres|cri\deleted{è}t\deleted{t}ères|cri\deleted{t}è\deleted{t}ères} signalés dans l'exemple 1 pour que le foncteur réalisation géométrique soit fidèle resp. pleinement fidèle, s'étendent de façon immédiate au cas envisagé ici.

\uline{Exemple} 3. Supposons que $\underline{R}=(\underline{C},\underline{s})$ satisfasse aux conditions suivantes. Le foncteur $\underline{s}$ se factorise en $\underline{C}\overset{\underline{s}'}{\qquad}\underline{\mathrm{F\unsure{\deleted{a}l|\deleted{l}a|\deleted{o}l|\deleted{l}o}op}}_S\qquad\underline{\mathrm{Fl}}_S$, où $\underline{\mathrm{Flop}}_S$ désigne la catégorie fibrée sur $\underline{\mathrm{Et}}_S$ des schémas relatifs à groupe fini (opérant à droite) d'opérateurs, et $\underline{\mathrm{Flop}}_S\quad\underline{\mathrm{Fl}}_S$ est le foncteur ``oubli\deleted{\unsure{r|e}} des opérations''. \deleted{Pour toute flèche $M'\quad M$ de $\underline{C}_{S'}$, la flèche \unsure{correspondante}}\inserted{Il faut préciser que (en vue de l'exemple kummérien plus} bas) nous prenons comme morphismes d'un schéma à opérateurs $(Z,G)$ sur $S$ dans un autre $(Z',G')$ les couples $(f,u)$, où $f:X\quad X'$ est un $S$-morphisme et $u:G_S\quad G'_S$ un morphisme(localement constant,\deleted{\unsure{(|f}}pas nécessairement constant !) de schémas en groupes sur $S$, de telle façon que $f$ soit compatible avec $u$ dans le s ens évident. On suppose que pour toute flèche $M'\quad M$ dans un $\underline{C}_S$, \deleted{déf} la flêche correspondante $\underline{s}(M')=(Z',G')\quad\underline{s}(M)=(Z,G)$ induit un isomorphisme $Z\quad Z'/H$, où $H$ est le noyau de $G'_S\quad G_S$ (qui est un sous-groupe localement constant de $G'_S$). On suppose \deleted{que tout} aussi que pour tout objet $M$ de $\mathrm{Ob}\,\underline{C}_{S'}$, posant $\underline{s}(M)=(Z,G)$, \deleted{et pour tout $g\quad G$, il existe un endomorphisme $\underline{g}$ de $M$, uniquement déterminé, tel que}\inserted{on se donne un \deleted{automorphisme} homomorphisme $G\quad\mathrm{Aut}(M)$, tel que pour tout $g\quad G$, si $\underline{g}$ est son image,} $\underline{s}'(\underline{g})$ soit le couple $(g_Z,\mathrm{int}(g))$, et \unsure{\inserted{tel} que tout|que tout \inserted{tel}} endomorphisme de $M$ soit localement de la forme $\underline{g}$; on \deleted{veut} suppose de plus que \deleted{la} l'homomorphisme $G\quad\mathrm{Aut}(M)$ soit compatib\unsure{\deleted{l}d|\deleted{d}l|\deleted{o}l|\deleted{l}o}e avec images inverses, et que pour toute flèche $M'\quad \unsure{M|\deleted{M}M|\deleted{X}M|\deleted{M}X}$ dans $\underline{C}_{S'}$, considérons $M'$ (resp. $M$) comme muni ud groupe $G'$ ($G$) d'opérateurs, $M'\quad M$ soit compatible avec $G'\quad G$.
% archive sheet number: 18
